\section{Conclusion}
\label{sec:conclusion}

HMRS connects the state used to manage collections, documents and segments
to the interfaces that agents use to find and inspect evidence. The
collection catalogue and document shelf expose metadata for narrowing
search, and opened documents and indexed segments carry the provenance
needed to inspect a result. Lifecycle checks make discrepancies observable:
a rename left duplicate vector entries, and re-indexing restored the expected
state.

On the 72-question evaluation, the tiered configuration that combined
document and chunk windows returned packs containing more annotated target
strings than the baseline configuration. The difference was 15 questions in the
frozen run and 12 in each rerun on the same questions.
In the frozen trace, the combined configuration used a higher recorded text
volume than the baseline and less than the mechanical full scan; the scan was
incomplete or absent in the later runs. In the historical verifier
comparison, keeping retained candidates in fusion order placed a gold
document in the top five for 15 of 16 queries, compared with seven under
verifier-score ordering.
These results motivate two concrete practices in managed document systems:
audit the state that supplies retrieval candidates, and test evidence
admission and ordering as separate decisions.
